% ===== 原文 PDF pp.55--58（印刷 pp.54--57）=====
% tr-05（2/3）。首行紧接 tr05a 末句（同一自然段，文件间不留空行）。
% [原文 p.54]
如果我们进一步把这样的理论紧化到 $\cC$ 上，那么 3d 理论的模空间就由 5d 场在 $\cC$ 上的 BPS 方程的解集所支配，而这些方程正是 $\cC$ 上的 Hitchin 方程组 [198]。如下文所解释的，HK 流形 $\cM$ 与 $\cC$ 上 Hitchin 方程的解模空间重合。

更具体地说，让我们描述由 6d $\cN=(2,0)$ 理论先在圆 $S^1$ 上、再在 $\cC$ 上紧化所得的 3d $\cN=4$ 理论的模空间。若记 5d SYM 理论的标量为 $\Phi_I$，$I=1,\ldots,5$，Hitchin 场为 (4.3) 中的组合，而规范场为
\begin{equation}
  A = A_z\dd z + A_{\bar z}\dd\bar z\,,
  \tag{4.6}
\end{equation}
那么 Hitchin 方程如下：
\begin{equation}
\begin{aligned}
  F + [\Phi_z,\bar\Phi_z] &= 0\,,\\
  \bar\partial_A\Phi_z &\equiv \dd\bar z\left(\partial_{\bar z}\Phi_z+[A_{\bar z},\Phi_z]\right) = 0\,,\\
  \partial_A\bar\Phi_z &\equiv \dd z\left(\partial_z\bar\Phi_z+[A_z,\bar\Phi_z]\right) = 0\,.
\end{aligned}
\tag{4.7}
\end{equation}
% [Tutor Note: (4.7) 的三条方程就是 5d $\cN=2$ SYM 的 BPS 方程（自对偶条件）在 $\cC$ 上的约化：$F$ 是 5d 规范场沿 $\cC$ 的分量的场强，第一条是 $F_{z\bar z}$ 的方程，后两条是 5d 标量的协变恒定性条件。$\bar\partial_A$、$\partial_A$ 分别表示协变导数的 $(0,1)$、$(1,0)$ 部分，故 $\bar\partial_A\Phi_z$ 是 $(0,1)$-形式、$\partial_A\bar\Phi_z$ 是 $(1,0)$-形式，各自前面正好配上 $\dd\bar z$ 与 $\dd z$。]
这些方程在 $\cC$ 上的解空间，描述了 3d 理论模空间中一个使扭曲紧化里的 $\fso(3)_R$ 对称性得以保持的分支（见 (4.2)）；这个分支应被认同为前面提到的 HK 流形 $\cM$，它也称为 Hitchin 模空间。

至此我们知道，用两种方式紧化 6d 理论所得的 3d 理论应当是同一个。然而，两个 RG 流之间存在一个非平凡的映射（见图 7）：如 [29] 所解释，它们由 3d 镜像对称相联系 [199]（另见 [200]）。\footnote{这里给出一个快速的论证。把 5d $\cN=2$ SYM 看作 $\cN=1$ 理论时，其超多重态中有两个复标量；(4.3) 中的标量 $\Phi_z$ 是其中之一，它在 $\fu(1)_r$ 下带荷，因而与描述 4d $\cN=2$ 的 CB 相关。因此，$\Phi_z$ 自然地既参数化约化到 $S^1$ 上的 4d $\cN=2$ 理论的 CB，又参数化约化到黎曼面 $\cC$ 上的 5d $\cN=2$ SYM 的 HB。} 若想了解 4d 理论的库仑分支，我们可以看 3d 理论的库仑分支，后者继承了前者，并因绕圆缠绕的线算子而得到增强。我们刚刚还说明了：保持 $\fso(3)_R$ 对称性的 3d 理论模空间，可以由 5d SYM 理论在黎曼面上紧化所得的 Hitchin 系统描述 [29,198,200]，其解参数化 5d 理论的希格斯分支。由 3d 镜像对称，由圆 $S^1$ 上紧化所得的 3d 库仑分支，与紧化在 $\cC$ 上的 5d 理论的希格斯分支相同，而后者正是 HK 流形 $\cM$。

% [原文 p.55]
最后，若我们把绕紧化圆的 4d 线算子上的坐标（它们参数化 $\cM$ 的纤维方向）抑制掉，4d 库仑分支就等同于 Hitchin 可积系统的底空间 $B$。举例来说，若 $\fg=A_{k-1}$，则底空间 $B$ 为
\begin{equation}
  B = \bigoplus_{r=2}^{k}H^0\!\left(\cC,\,K_\cC^{\otimes r}\right),
  \tag{4.8}
\end{equation}
它们是形如 $\Tr(\Phi_z^r)$ 的 Hitchin 场规范不变单项式的全纯截面。于是 Seiberg--Witten 曲线由可积系统 (4.7) 的谱曲线给出，而它恰好就是 (4.4)。
% [Tutor Note: (4.8) 的求和从 $r=2$ 起：$\fg=A_{k-1}=\fsu(k)$ 的 $r=1$ 项是 $\Tr\Phi_z$，它在 $\fsu(k)$ 中恒为零（无迹），所以 Hitchin 底空间不含 $r=1$ 的截面；于是求和恰有 $k-1$ 项，正好等于 $\fsu(k)$ 的独立 Casimir 数（= 秩 $k-1$），库仑分支维数随之等于 $k-1$。$K_\cC$ 是 $\cC$ 上的典范线丛，$K_\cC^{\otimes r}$ 的全纯截面就是 $r$ 次全纯微分。]

让我们最后一次强调：我们是从描写 5d $\cN=2$ SYM 模空间的 BPS 方程的经典解，得到了一个一般而言强耦合的 4d $\cN=2$ 理论的量子库仑分支 EFT。这之所以可行，是因为镜像对称把前者与镜像理论的希格斯分支联系起来，因而前者受到保护、不受量子修正影响 [201]。
% [Tutor Note: 这一步的合法性全在「镜像对称」四个字上：Hitchin 模空间是 5d SYM 在 $\cC$ 上的经典（BPS）解空间，而它与 3d 理论保持 $\fso(3)_R$ 的分支之间的等同是由镜像对称保证的；后者是 BPS 保护的量，故不受量子修正。这就是「由经典 Hitchin 系统得到量子几何」的确切含义——不是说经典解本身精确，而是说镜像对偶把经典解映射到了 3d 的精确分支上。]

\SUB{4.5}{$\cC$ 上的余维二缺陷与穿刺点}

正如我们在上一节所见，4d $\cN=2$ 理论的可积系统已与 $\cC$ 上规范群为 $\SU(N)$ 的 Hitchin 系统相联系 [197]。我们把 $\Phi_z$ 视为 Hitchin 系统的 Higgs 场。它是 $\SU(N)$ 伴随表示中的全纯 1-形式。该系统的 SW 曲线 $\Sigma$ 由式 (4.4) 给出。通过计算行列式，我们可以把 SW 曲线改写为
\begin{equation}
  x^n = \sum_{i=2}^{n}\phi_i(z)\,x^{n-i}\,,
  \tag{4.9}
\end{equation}
其中 $\phi_i(z)$ 是定义在带穿刺点黎曼面上的 $i$ 次全纯函数。对规范群为 $\SU(N)$ 的情形，$\phi_j$ 是 Higgs 场规范不变组合即 $\Tr(\Phi_z^j)$ 的多项式。用 Hitchin 系统的语言说，$\cC$ 上的点状缺陷对应于 $\Phi_z$ 的一个极点（也就是相应的边界条件）。若取局部坐标 $z$ 使穿刺点位于 $z=0$，则 Higgs 场在局部表现为
\begin{equation}
  \Phi_z = \frac{A}{z} + \ldots
  \tag{4.10}
\end{equation}
% [Tutor Note: 由 (4.4) 展开 $\det(x-\varphi_z)=0$ 即得 (4.9)：$\varphi_z$ 的 $n$ 个本征值给出 $x$，而 $\det$ 展开的系数正是 $\phi_i=\Tr(\varphi_z^i)$（差一些对称化因子）。这也说明 (4.9) 与 (4.4) 是同一个方程，只是把行列式写成特征多项式的形式；因此 (4.10) 中 $A$ 的极点会原样传给每一个 $\phi_i$，极点阶数由 $A$ 的 Jordan 结构决定（见 4.5 节后半）。]
% [原文 p.56]
其中 $A$ 是 $\fsu(N)$ 的一个元素，它指定穿刺点的性质，并影响由 6d 自对偶弦沿围绕极点的闭路缠绕所得的 BPS 粒子谱。由于 $\Phi_z$ 不是规范不变的，这些缺陷由 $A$ 的共轭类刻画，等价地由 $\fsu(N)$ 中的（余）伴随轨道刻画（综述见 [41]）。这类轨道有两类，对应于矩阵 $A$ 的 Jordan 型：半单矩阵给出所谓的 \textbf{半单}轨道，幂零矩阵给出 \textbf{幂零}轨道。在前一种情形，我们引入了额外的标度，即 $A$ 的本征值，它们在 $\fu(1)_r$ R 对称性下带荷。这些本征值由 SW 微分绕穿刺点的积分拾取，并编码低能 EFT 中 BPS 粒子的质量。因此，我们把 $A$ 的本征值称为穿刺点的质量参数；它们自然地破坏共形不变性。另一方面，幂零的 $A$ 不引入任何物理标度，因为由 $A$ 的幂次构造出的任何规范不变组合（这类组合本会引入这样的标度）都为零。\footnote{等价地，对幂零穿刺点作一次标度变换，总可以借助一次复化规范变换把它抵消掉。} 因此，以幂零轨道标记的穿刺点有助于构造明显保持共形对称性的理论。

\begin{examplebox}{例 4.1：}
  让我们以 $\fsu(2)$ 规范理论为例来说明。在这种情况下我们只有一个 Casimir 算子
  \begin{equation}
    \phi_2 = \frac{1}{2}\Tr\left(\Phi_z^2\right),
    \tag{4.11}
  \end{equation}
  其中
  \begin{equation}
    \Phi_z \sim \frac{1}{z}
    \begin{pmatrix}
      m & 0\\
      0 & -m
    \end{pmatrix}
    + \text{正则项}\,.
    \tag{4.12}
  \end{equation}
  相应的 Casimir 于是具有如下形式
  \begin{equation}
    \phi_2(z) = \frac{m^2}{z^2} + \ldots
    \tag{4.13}
  \end{equation}
  其中 $m$ 是一个\emph{质量参数}。该参数代表把 Seiberg--Witten 微分绕 $z=0$ 积分时出现的 BPS 粒子的质量。次领头项编码 CB 算子的 VEV。

  % [原文 p.57]
  在无质量极限下，极点仍是一阶的，
  \begin{equation}
    \Phi_z \sim \frac{1}{z}
    \begin{pmatrix}
      0 & 1\\
      0 & 0
    \end{pmatrix}
    + \text{正则项}\,,
    \tag{4.14}
  \end{equation}
  $\phi_2$ 也类似。
\end{examplebox}

% [Tutor Note: (4.12) 的极点阶数是 $1$，但 $m\ne 0$ 时 $\phi_2=\tfrac12\Tr\Phi_z^2$ 的极点阶数是 $2$，这正对应 $\{p_2\}=2$：$A$ 的极点阶 $p_j$ 与 $\phi_j$ 的极点阶一致。$m\to 0$ 时 $\Phi_z$ 变成幂零矩阵 $\begin{pmatrix}0&1\\0&0\end{pmatrix}$（Jordan 块 $J_2$），极点阶不变但 $\phi_2$ 的极点阶降为 $1$，这就是「正则穿刺点里最简单的极小穿刺点」的由来。]

\begin{exercisebox}{习题 4.1}
  证明：借助复化规范群 $\SL(2,\C)$ 中的矩阵作共轭，可以把 $\begin{pmatrix} m & 0\\ 0 & -m\end{pmatrix}$ 化为 $\begin{pmatrix} m & a\\ 0 & -m\end{pmatrix}$ 的形式，其中 $a\in\C$。
\end{exercisebox}
% [Tutor Note: 习题 4.1 的要点：$\begin{pmatrix} m & a\\ 0 & -m\end{pmatrix}$ 的本征值仍是 $\pm m$，且因 $m\ne 0$ 时两本征值不相等，它是可对角化的，故与 $\diag(m,-m)$ 共轭；$\SL(2,\C)$ 中的矩阵可以连续地把非对角元 $a$ 调到任意值。物理上这说明第一类（半单）穿刺点的质量参数只由本征值决定，与 $a$ 无关。]

\begin{center}
  \includegraphics[width=0.42\linewidth]{fig8.png}
\end{center}
\figtitle{8}{$T_2$ 理论。}

更具体地说，让我们考虑由带 3 个穿刺点的球面 $S^2$ 定义的 UV 曲线 $\cC$，如式 (4.12) 那样。这有时被称为 $T_2$ 理论，如图 8 所示。我们用 $m_1$、$m_2$ 与 $m_3$ 分别表示位于 $z=0$、$1$、$\infty$ 的穿刺点的质量参数。Casimir $\phi_2(z)$ 可写成
\begin{equation}
  \phi_2(z) = \frac{f(z,m_i)}{z^2(z-1)^2}\,,
  \tag{4.15}
\end{equation}
其中 $f(z,m_i)$ 是 2 次多项式，使得 Hitchin 场在 $0$、$1$、$\infty$ 处的穿刺点附近表现为式 (4.12) 那样。该理论的库仑分支是平凡的，EFT 除质量 $m_i$ 之外不含任何其他参数，正是由于这个原因它有时被称为「刚性」的。$T_2$ 理论确实描述了一个自由的 $\cN=2$ 半超多重态，它处于 $\SU(2)^3$ 的三重基本表示中。
% [Tutor Note: (4.15) 的分母 $z^2(z-1)^2$ 保证 $\phi_2$ 在 $z=0$ 与 $z=1$ 各有二阶极点（对应 (4.12) 中 $\Phi_z$ 的一阶极点）；分子取 2 次多项式，则用 $w=1/z$ 代入即可看出 $z\to\infty$ 处同样只有二阶极点。$f$ 的 3 个系数恰好被三个条件 $f(0)=m_1^2$、$f(1)=m_2^2$ 与 $z\to\infty$ 的首项系数 $=m_3^2$ 完全定死，没有剩余自由度——这正是该理论库仑分支平凡（「刚性」）的原因。]

这里我们聚焦于 \textbf{正则}（或称 \textbf{驯服}）穿刺点，即 Hitchin 场在穿刺点位置处领头阶为式 (4.12) 那种单极点的情形。
